% ================= SINGLE-SOURCED DESIGN VALUES =================
% Every recurring numerical value of the converged design point is defined
% ONCE here and referenced everywhere else, so the same quantity cannot carry
% two different figures in two different tables. This is the document-side
% counterpart of CAD convention 1 (Section~\ref{sec:cad_method}): no table
% hard-codes a value that duplicates one of these.
%
% All of these are regenerated by analysis/SIZING.py. If that run changes,
% change the macro here and the whole document follows; do not edit the
% number inside a table.
%
% Usage: \qty{\Mbasis}{\kilo\gram}, \qty{\wheelOD}{\milli\metre}, ...
% Values are bare numbers (no units, no \qty) so they compose with siunitx.

% ---------- cube ----------
\newcommand{\cubeL}{150}          % [mm] outer edge length, FROZEN (CON-01, gate M2a)
\newcommand{\cubeLan}{149}        % [mm] edge length the sizing analysis is run at
\newcommand{\cubeWall}{5}         % [mm] frame wall thickness
\newcommand{\Mbasis}{1.800}       % [kg] sizing basis -- the mass the wheel was sized against
\newcommand{\Mcap}{1.8}           % [kg] CON-02 mass cap
\newcommand{\Mbuilt}{1.626}       % [kg] as-built, weighed (App.~\ref{app:cad_mass})
\newcommand{\Mdelta}{174}         % [g]  Mbasis - Mbuilt, realised margin
\newcommand{\MbasisG}{1800.0}     % [g]  sizing basis, in grams (mass-budget tables)

% ---------- mass budget ----------
\newcommand{\massKnown}{1367.5}   % [g] specified hardware, measured or datasheet
\newcommand{\massStructAllow}{432.5}  % [g] structure allowance, pending CAD closure
\newcommand{\massStructBuilt}{258.5}  % [g] structure implied by the weighed article
\newcommand{\massWheelsThree}{702.4}  % [g] three reaction wheels

% ---------- reaction wheel ----------
\newcommand{\wheelOD}{120}        % [mm] ring outer diameter, fixed by the envelope
\newcommand{\wheelID}{100}        % [mm] ring inner diameter
\newcommand{\wheelWr}{10}         % [mm] ring radial width = ballast hole depth
\newcommand{\wheelT}{20}          % [mm] wheel axial thickness
\newcommand{\wheelSpokes}{3}      % [-]  spoke count
\newcommand{\wheelSpokeW}{10}     % [mm] spoke width
\newcommand{\wheelRmean}{55.0}    % [mm] ballast centroid radius, forced by the ID-to-OD hole
\newcommand{\wheelHole}{6.4}      % [mm] M6 clearance hole diameter
\newcommand{\wheelN}{15}          % [-]  ballast stations populated
\newcommand{\wheelNmax}{37}       % [-]  stations the 2 mm gap rule admits
\newcommand{\wheelMass}{234.14}   % [g]  wheel mass as ballasted
\newcommand{\wheelPlastic}{121.63}% [g]  PET-CF in the finished wheel
\newcommand{\wheelSteel}{112.50}  % [g]  M6 steel in the finished wheel
\newcommand{\wheelMassSolid}{127.86}  % [g] unballasted ring + spokes
\newcommand{\wheelI}{6.2715e-4}   % [kg m^2] delivered inertia
\newcommand{\wheelIsolid}{3.0451e-4}  % [kg m^2] unballasted ring + spokes
\newcommand{\wheelIstation}{2.151e-5} % [kg m^2] net inertia added per station
\newcommand{\wheelMstation}{7.085}    % [g] net mass added per station
\newcommand{\wheelForce}{72}      % [N]  centrifugal load per station at the design speed

% ---------- jump-up operating point ----------
\newcommand{\rpmDesign}{4000}     % [rpm] design wheel speed
\newcommand{\rpmFloor}{3473}      % [rpm] lowest speed the built wheel still closes at, as built
\newcommand{\rpmCeilPc}{57.8}     % [%]  design speed as a fraction of the derived ceiling (PER-08)
\newcommand{\hwDesign}{0.2612}    % [kg m^2/s] required per-wheel momentum at the sizing basis
\newcommand{\Itarget}{6.2352e-4}  % [kg m^2] inertia requirement at the sizing basis (PER-06)
\newcommand{\ItargetBuilt}{5.4445e-4} % [kg m^2] the same requirement at the as-built mass
\newcommand{\marginSized}{100.6}  % [%]  delivered inertia against the sizing basis
\newcommand{\marginBuilt}{115.2}  % [%]  delivered inertia against the as-built mass
\newcommand{\tauBrake}{5}         % [N m] brake torque per wheel, sizing value
\newcommand{\tauGrav}{1.3155}     % [N m] gravitational tip-over torque at the sizing basis
\newcommand{\tauGravBuilt}{1.1884}% [N m] the same at the as-built mass
\newcommand{\betaBrake}{1.3570}   % [-]  brake-amplification factor at the sizing basis
\newcommand{\tspin}{1.2}          % [s]  spin-up to design speed at the continuous current limit
\newcommand{\tspinFloor}{0.20}    % [s]  spin-up floor set by the tip-over torque
\newcommand{\wheelKE}{55}         % [J]  kinetic energy stored per wheel at design speed

% ---------- actuator ----------
% The design torque constant is the MACHINE constant implied by the datasheet
% KV rating, k_t = 60/(2*pi*KV) = 60/(2*pi*380) = 0.02513 N m/A. It is a
% property of the motor and independent of what the shaft drives, which is why
% it -- and not a fit to the propeller bench data of Section 5.4.1 -- is the
% figure the design runs on. That fit is a cross-check only: see the discussion
% there of why a straight line through a curved torque-current characteristic
% returns a slope that understates torque per ampere inside the operating band.
% Rig measurement (WBS E6) replaces this figure; until then every torque/current
% pair in the document must be computed from \ktMotor and nothing else.
\newcommand{\motorKV}{380}        % [rpm/V] MN4006 KV rating (Tab. tab:bom-supplied)
\newcommand{\ktMotor}{0.02513}    % [N m/A] motor torque constant, from \motorKV
\newcommand{\tauDriveCont}{0.226} % [N m] drive torque at the 9 A driver continuous limit
\newcommand{\tauDrivePeak}{0.402} % [N m] drive torque at the 16 A motor short-duration limit
\newcommand{\pcDriveGrav}{17.2}   % [%]  \tauDriveCont as a fraction of \tauGrav
